\chapter*{About This Handout}\label{chap:about}
\addcontentsline{toc}{chapter}{About This Handout}

This handout goes with the \hologo{LaTeX} workshop.
It contains everything the slides show,
plus the details the slides left out on purpose.
It is meant as a reference for later, for your next lab report and for your thesis,
so you can also read it without having attended the workshop.

\section*{How it is organised}

\autoref{chap:introduction} explains what \hologo{LaTeX} is,
and \autoref{chap:setup} how to install it or use it online.
Chapters~\ref{chap:skeleton} to~\ref{chap:references} follow the four blocks of the workshop:
the skeleton of a document, text, math, and then tables, figures, cross-references and citations.
Chapters~\ref{chap:errors} to~\ref{chap:resources} go beyond the slides:
error messages, useful packages, larger documents and where to look things up.
The appendix holds the complete example document and a cheat sheet.

\section*{The running example}

Most examples come from a short lab report on Newton's law of cooling,
\enquote{How Fast Does Your Coffee Go Cold?}.
It compares how fast a drink cools in a mug, a glass and a vacuum flask.
In the workshop, you write this report yourself from a template,
step by step from an empty document to a report with equations, a table, a figure and a bibliography.
The template, with its finished version, can be downloaded from the repository (see \enquote{\nameref{sec:about-source}} below).

\section*{Code and output}

Code is shown in shaded boxes like this one:

\begin{codeblock}
The underlying physics is \textit{Newton's law of cooling}.
\end{codeblock}

The small grey numbers in the margin are line numbers.
They are not part of the code; the text uses them to point at a line.

What \hologo{LaTeX} makes of the code is shown in a box with a thin frame underneath:

\sliderender{01-styling.pdf}

These outputs were compiled separately, as a real document
with the class \texttt{article} (\texttt{a4paper}, \texttt{11pt})
and the packages the workshop has introduced up to that point.
So they look exactly like your own document, not like this handout,
which uses a different font (Times) and loads many more packages.

Most examples are only a piece of a document.
Unless the box says otherwise,
the code belongs between \code|\begin{document}| and \code|\end{document}|.
Code with the title \enquote{Preamble} goes before \code|\begin{document}|
(\autoref{sec:preamble} explains what the preamble is):

\begin{codeblocktitled}{Preamble}
\usepackage{parskip}
\end{codeblocktitled}

\section*{Conventions}

\begin{itemize}
    \item Commands, file names and anything you type literally are set in \texttt{typewriter} font.
    \item File names are always written with their extension (\texttt{.tex}, \texttt{.bib}, \texttt{.pdf}),
        so it is clear which file is meant.
    \item Links to websites are \myhref{https://www.overleaf.com}{blue and underlined},
        and the link text is usually the address itself, so it can be typed from a printed copy.
        References inside the handout (\enquote{\autoref{chap:setup}}) are black, but just as clickable.
\end{itemize}

\section*{The source of this handout}\label{sec:about-source}

This handout is written in \hologo{LaTeX} itself,
and it uses most of what it describes:
\texttt{booktabs} tables, \texttt{siunitx} units, \texttt{\textbackslash autoref}, \texttt{biblatex}.
Its source, the slides and the workshop template are in the repository
\myhref{\repourl}{\repodisplay};
the template's \texttt{.zip} file for Overleaf is attached to its releases.
